\section{Key Stakeholders \& Blocs}

\subsection{Western Enforcement Powers (France, Germany, United Kingdom, United States, Canada, Australia)}

These delegations represent the principal architects of the contemporary international anti-money laundering framework. As G7 members and founding drivers of the Financial Action Task Force, they have shaped the global regulatory architecture governing illicit finance for over three decades, from the 1989 FATF mandate through the post-2008 push for automatic exchange of financial information and the more recent beneficial ownership transparency agenda.\footnote{Financial Action Task Force. (2023). \emph{Revised FATF recommendations}. \url{https://www.fatf-gafi.org/}} Their position in this committee will consistently favor expanding INTERPOL's operational capacity, strengthening cross-border information sharing obligations, and imposing stricter due diligence requirements on professional enablers such as lawyers, accountants, and trust administrators.\footnote{INTERPOL. (2023). \emph{Financial crimes strategy}. \url{https://www.interpol.int/}} \footnote{Nance, M. T. (2018). The regime that FATF built: An introduction to the Financial Action Task Force. \emph{Crime, Law and Social Change, 69}(2), 109–129.}

The United States occupies a particularly dominant role through FinCEN, the Foreign Corrupt Practices Act, and the extraterritorial reach of dollar-denominated correspondent banking.\footnote{Farrell, H., \& Newman, A. L. (2019). Weaponized interdependence: How global economic networks shape state coercion. \emph{International Security, 44}(1), 42–79.} The 2021 Anti-Money Laundering Act and the Corporate Transparency Act represent the most significant domestic reform efforts in decades, though implementation remains contested.\footnote{Cassella, S. D. (2022). The Anti-Money Laundering Act of 2020: Congress addresses the deficiencies in U.S. AML law. \emph{Journal of Money Laundering Control, 25}(3), 473–488.} The United Kingdom has positioned itself as a transparency champion through its public beneficial ownership register while simultaneously facing criticism for London's role in concentrating post-Soviet oligarchic wealth in British real estate and corporate structures.\footnote{Bullough, O. (2022). \emph{Butler to the world: How Britain became the servant of tycoons, tax dodgers, kleptocrats and criminals}. Profile Books.} Delegates should note that European delegations tend to favor multilateral rule-based solutions, while the United States has historically preferred bilateral enforcement and unilateral extraterritorial action, a tension that will surface in debates over whether new INTERPOL mechanisms should be binding or recommendatory.\footnote{Zarate, J. C. (2013). \emph{Treasury's war: The unleashing of a new era of financial warfare}. PublicAffairs.}

\subsection{Offshore Financial Centers (Switzerland, Luxembourg, Panama, Bahamas, Monaco, Liechtenstein)}

This bloc occupies the most structurally complex and politically sensitive position in the committee. These states share a common economic interest in preserving the conditions of legal confidentiality, low taxation, and flexible corporate law that attract foreign capital, but differ significantly in their degree of formal compliance with international norms. \footnote{Palan, R., Murphy, R., \& Chavagneux, C. (2010). \emph{Tax havens: How globalization really works}. Cornell University Press.}

Switzerland and Luxembourg represent the institutionally mature end of the spectrum. Both have undergone substantial reform since the 2008 financial crisis, adopting automatic exchange of information under the Common Reporting Standard and nominally complying with FATF recommendations.\footnote{Organisation for Economic Co-operation and Development. (2017). \emph{Standard for automatic exchange of financial information in tax matters: Implementation handbook} (2nd ed.). OECD Publishing.} Their delegates will argue that secrecy has been replaced by privacy, a legally calibrated protection for legitimate wealth rather than a shield for criminality.\footnote{Emmenegger, P. (2017). Swiss banking secrecy and the problem of international cooperation in tax matters: A nut too hard to crack? \emph{Regulation \& Governance, 11}(1), 24–40.} This position is contested by studies of the Tax Justice Network that have consistently ranked Switzerland among the world's largest enablers of illicit financial flows even after its formal reforms, largely because of the gap between legal compliance and operational transparency. Panama occupies a uniquely exposed position following the 2016 Panama Papers, which revealed how the law firm Mossack Fonseca helped clients across 200 countries conceal assets through anonymous shell companies.\footnote{Obermayer, B., \& Obermaier, F. (2016). \emph{The Panama Papers: Breaking the story of how the rich and powerful hide their money}. Oneworld Publications.} Caribbean offshore centers such as the Bahamas will argue that enforcement frameworks designed by and for large wealthy economies disproportionately punish small island economies with limited alternative development pathways.\footnote{Vlcek, W. (2017). \emph{Offshore finance and small states: Sovereignty, size, and money}. Palgrave Macmillan.}

\subsection{Emerging Economies (China, India, Brazil, South Africa)}

The emerging economies bloc approaches this topic from a structurally ambivalent position. These states are among the most significant victims of illicit financial flows globally. The African Union estimates that Africa alone loses over \$50 billion annually to illicit outflows, while Global Financial Integrity has documented comparable losses across Latin America and South Asia.\footnote{United Nations Economic Commission for Africa \& African Union. (2015). \emph{Report of the High Level Panel on Illicit Financial Flows from Africa}. \url{https://www.uneca.org/}} \footnote{Global Financial Integrity. (2021). \emph{Trade-related illicit financial flows in 135 developing countries: 2009–2018}. \url{https://gfintegrity.org/}} Yet these same states have historically been among the most vocal critics of the international enforcement architecture, arguing that FATF's listing and delisting processes reflect the geopolitical priorities of its predominantly Western membership rather than objective technical assessment of financial crime risk.\footnote{Nance, M. T. (2018). The regime that FATF built: An introduction to the Financial Action Task Force. \emph{Crime, Law and Social Change, 69}(2), 109–129.} \footnote{Soederberg, S. (2004). \emph{The politics of the new international financial architecture: Reimposing neoliberal domination in the Global South}. Zed Books.} Countries across the Global South have been grey or blacklisted for regulatory gaps while jurisdictions such as the United States, which did not implement the Common Reporting Standard, face no comparable scrutiny. This bloc will be broadly supportive of transparency in principle but insist that any new framework be accompanied by governance reform giving emerging economies a meaningful voice in standard-setting.\footnote{Vestergaard, J., \& Wade, R. H. (2012). The governance response to the great recession: The "success" of the G20. \emph{Journal of Economic Policy Reform, 15}(2), 71–86.} China will support INTERPOL's operational role selectively, where it aligns with domestic enforcement priorities, while opposing frameworks that could subject Chinese sovereign financial activity to external oversight.\footnote{Liss, C., \& Sharman, J. C. (2015). Global corporate crime-fighters: Private transnational responses to piracy and money laundering. \emph{Review of International Political Economy, 22}(4), 693–718.}

\subsection{Gulf States (UAE, Saudi Arabia, Qatar)}

The Gulf delegations bring sovereign wealth management concerns, regional geopolitical calculations, and ongoing credibility questions to this committee. The UAE, and Dubai specifically, has become one of the world's most significant centers for illicit financial flows, functioning as a destination for Russian oligarchic capital following 2022 sanctions, Iranian sanctions evasion, and South Asian hawala networks.\footnote{Carnegie Endowment for International Peace. (2020). \emph{Dubai's role in facilitating corruption and global illicit financial flows}. \url{https://carnegieendowment.org/}} FATF placed the UAE on its grey list in 2022, citing deficiencies in oversight of real estate transactions, precious metals trading, and corporate service providers.\footnote{Financial Action Task Force. (2022). \emph{Jurisdictions under increased monitoring: United Arab Emirates}. \url{https://www.fatf-gafi.org/}} The UAE was removed from the grey list in 2024 following formal reform commitments, though independent analysts have questioned the depth of implementation.\footnote{Chatham House. (2024). \emph{Gulf states and financial regulation: Progress and persistent gaps}. The Royal Institute of International Affairs.} Gulf delegations will support INTERPOL cooperation on terrorism financing where security interests align with Western partners, while drawing clear limits around provisions that would extend international oversight to sovereign capital flows or real estate markets.\footnote{Hertog, S. (2010). \emph{Princes, brokers, and bureaucrats: Oil and the state in Saudi Arabia}. Cornell University Press.}

\subsection{Geopolitically Isolated States (Russia, Iran)}

These two delegations occupy a categorically distinct position in this committee as the primary targets of the most extensive international financial sanctions regimes currently in operation.\footnote{Hufbauer, G. C., Schott, J. J., Elliott, K. A., \& Oegg, B. (2007). \emph{Economic sanctions reconsidered} (3rd ed.). Peterson Institute for International Economics.} For Russia, the freezing of approximately \$300 billion in sovereign assets following the 2022 invasion of Ukraine represents what Moscow frames as the weaponization of international financial infrastructure for geopolitical ends, a precedent it will argue undermines the legitimacy of the entire framework within which INTERPOL operates.\footnote{Connolly, R. (2018). \emph{Russia's response to sanctions: How Western economic statecraft is reshaping political economy in Russia}. Cambridge University Press.} Russia will push for language explicitly protecting states from politically motivated financial investigations.\footnote{Drezner, D. W. (2022). \emph{The uses and abuses of weaponized interdependence}. Brookings Institution Press.} Iran, facing near-total exclusion from formal financial channels, has driven reliance on hawala, gold-based transfers, and cryptocurrency.\footnote{Maloney, S. (2015). \emph{Iran's political economy since the revolution}. Cambridge University Press.} Both delegations will contest the neutrality of INTERPOL's mechanisms and will be essential voices in any debate over the boundaries of INTERPOL's mandate and the protection of sovereignty within international financial governance.
